% Artigo científico — submissão PUCRS (minuta para revisão do autor)
% Roteamento por atenção com garantias verificadas por máquina
\documentclass[12pt,a4paper]{article}
\usepackage[utf8]{inputenc}
\usepackage[T1]{fontenc}
\usepackage[brazilian]{babel}
\usepackage{amsmath,amssymb}
\usepackage{booktabs}
\usepackage{hyperref}
\usepackage{geometry}
\geometry{margin=3cm}
\title{Roteamento por Atenção com Garantias Verificadas por Máquina:\\
Teoremas em Lean 4, Calibração Empírica Interna e Limites Declarados}
\author{Marcelo Claro Laranjeira \\ Mestre em Ciência da Educação \\ Pontifícia Universidade Católica do Rio Grande do Sul (PUCRS) \\ ORCID: \url{https://orcid.org/0000-0001-8996-2887} \\ \texttt{marceloclaro@gmail.com} --- \url{https://bit.ly/geomaker}}
\date{Outubro de 2026 --- minuta, não submetida}
\begin{document}
\maketitle

\begin{abstract}
\textbf{Resumo.} Este artigo relata a formalização e a calibração interna de um
mecanismo de roteamento por atenção para orquestração de agentes, inspirado no
mecanismo de atenção do Transformer. Onze lemas foram verificados por máquina
(\textit{machine-checked}) em Lean 4 (v4.34.1) com Mathlib (rev.\ v4.34.1):
normalização do softmax, limitação bilateral da combinação convexa, existência
de candidato de utilidade máxima, invariância à subtração do máximo, pesos em
$(0,1)$ e peso unitário no caso singleton, monotonicidade estrita (T8),
dominação top-\textit{k} (T9), nitidez por temperatura (T10) e ausência de
overflow (T11); todos dependem apenas dos axiomas
padrão (\textit{propext}, \textit{Classical.choice}, \textit{Quot.sound}), sem
\textit{sorry}. Uma bancada sintética interna (200 decisões, 5 sementes)
comparou o roteador real a baselines ingênuos: arrependimento (\textit{regret})
médio de $0{,}042\pm0{,}011$ contra $0{,}19$--$0{,}22$ dos baselines, com
intervalos de confiança sem sobreposição. Ablações e análise de sensibilidade
delimitam a validade dos achados, incluindo a dependência da cabeça de
confiança. Nada aqui autoriza alegações de superioridade externa, inovação ou
disrupção: trata-se de evidência preliminar com trilha de reprodução publicada.
\\
\textbf{Palavras-chave:} assistentes de prova; Lean 4; Mathlib; roteamento por
atenção; verificação formal; reprodutibilidade.
\end{abstract}

\begin{abstract}
\textbf{Abstract.} This paper reports the formalization and internal
calibration of an attention-based routing mechanism for agent orchestration,
inspired by the Transformer attention mechanism. Eleven lemmas were
machine-checked in Lean 4 (v4.34.1) with Mathlib (rev.\ v4.34.1): softmax
normalization, two-sided convex-combination bounds, existence of a
utility-maximal candidate, max-subtraction invariance, weights in $(0,1)$,
and unit weight in the singleton case, strict monotonicity (T8), top-\textit{k}
dominance (T9), temperature sharpening (T10) and overflow-freedom (T11); all depend only on the standard axioms
(\textit{propext}, \textit{Classical.choice}, \textit{Quot.sound}), with no
\textit{sorry}. An internal synthetic benchmark (200 decisions, 5 seeds)
compared the real router against naive baselines: mean regret
$0.042\pm0.011$ versus $0.19$--$0.22$ for baselines, with non-overlapping
confidence intervals. Ablations and sensitivity analysis bound the findings,
including dependence on the confidence head. Nothing here supports claims of
external superiority, innovativeness, or disruption: this is preliminary
evidence with a published reproduction trail.
\\
\textbf{Keywords:} proof assistants; Lean 4; Mathlib; attention routing;
formal verification; reproducibility.
\end{abstract}

\section{Introdução}

A orquestração de agentes de software exige decidir, por tarefa, qual agente
executar sob critérios múltiplos (aderência semântica, capacidades,
confiança, carga). Mecanismos inspirados na atenção do Transformer
\cite{vaswani2017} ponderam candidatos por pesos normalizados, mas raramente
acompanham garantias formais sobre esses pesos ou sobre a otimalidade da
seleção. Este trabalho registra, com trilha auditável completa: (i) sete
lemas verificados por máquina sobre o núcleo matemático desse roteamento, em
Lean 4 com Mathlib \cite{moura2021,mathlib2020}; (ii) calibração empírica
interna contra baselines ingênuos; (iii) ablações, contraprovas numéricas e
revisões independentes internas; (iv) limites explícitos do que foi e do que
não foi demonstrado. A contribuição reivindicada é deliberadamente restrita:
otimalidade \textit{relativa} à utilidade definida, com reprodução publicada.

\section{Fundamentação}

O roteador estudado combina quatro cabeças normalizadas em $[0,1]$
(semântica $0{,}30$, capacidade $0{,}35$, confiança $0{,}25$, carga $0{,}10$),
exige soma unitária dos pesos sob pena de erro, aplica softmax com subtração
do máximo e ordena os candidatos do melhor ao pior. A Equação~\ref{eq:soft}
define os pesos; a Equação~\ref{eq:util} a utilidade combinada.
\begin{equation}
w_i = \frac{e^{f_i}}{\sum_j e^{f_j}} \label{eq:soft}
\end{equation}
\begin{equation}
u(a) = \sum_h \alpha_h \cdot s_h(a), \quad \sum_h \alpha_h = 1,\ \alpha_h \geq 0 \label{eq:util}
\end{equation}
Os onze lemas formalizam: soma unitária (T1), limitação bilateral por
combinação convexa (T2, T6), existência de maximizador em conjunto finito
não vazio (T3), invariância à subtração de constante (T4), pertinência a
$(0,1)$ com ao menos dois candidatos (T5), peso unitário no singleton (T7),
monotonicidade estrita no próprio escore (T8), dominação top-\textit{k} (T9),
nitidez por temperatura (T10) e ausência de overflow (T11).

\subsection{Esboço teórico: atenção como ponderação exponencial calibrada}

Dada uma lista finita de escores $f_1,\dots,f_n$, a Equação~\ref{eq:soft}
define a ponderação exponencial normalizada. Três fatos elementares,
aqui verificados por máquina, sustentam seu uso decisório. Primeiro, a massa
total é exatamente 1 (T1): os pesos são uma distribuição de probabilidade
sobre os candidatos, e a calibração empírica confirma soma 1 até
$1{,}1\times10^{-16}$. Segundo, deslocar todos os escores por uma constante
não altera nenhum peso (T4) --- é o que autoriza subtrair o máximo por
estabilidade numérica sem mudar a decisão; a contraprova em ponto flutuante
registrou invariância exata em 20\,000 vetores. Terceiro, cada peso está em
$(0,1)$ com ao menos dois candidatos (T5) e vale 1 no singleton (T7): nenhum
candidato é anulado nem absolutizado, salvo unicidade.

A correspondência com o Transformer é estrutural e parcial, não transferência
de propriedades: consulta $\leftrightarrow$ vetor da tarefa, chave
$\leftrightarrow$ vetor do cartão do agente, escore $\leftrightarrow$
similaridade e cabeças, valor $\leftrightarrow$ (utilidade, identificador),
agregação $\leftrightarrow$ combinação convexa seguida de softmax, decisão
$\leftrightarrow$ máximo. Nenhuma capacidade representacional ou desempenho
de Transformers treinados é herdada por essa analogia.

\subsection{Esboço matemático: combinação convexa e seleção gulosa}

Com pesos não negativos somando 1, a utilidade combinada da
Equação~\ref{eq:util} situa-se entre o mínimo e o máximo dos componentes
(T2, T6): o escore final nunca escapa ao intervalo das cabeças, o que dá
sentido formal ao validador que rejeita pesos cuja soma difere de 1. Em
conjunto finito não vazio existe maximizador (T3), e a regra top-1 o
seleciona; o padrão top-\textit{k} domina o resto sob ordem pairwise (T9), e
o peso cresce estritamente com o próprio escore (T8), o que justifica decidir
pelo máximo; a nitidez cresce ao esfriar (T10) e os expoentes deslocados não
estouram (T11); o custo documentado é o da ordenação completa onde bastaria
seleção parcial. Em Lean, as provas usam positividade da exponencial,
reescrita de somas (\texttt{Finset.sum\_div}, \texttt{Finset.sum\_mul}),
monotonicidade da multiplicação por não negativos e existência de máximo de
imagem finita --- todos sobre os axiomas padrão, auditados um a um.

\section{Método}

\textit{Provas.} Arquivo \texttt{roteamento\_atencao.lean} verificado por
\texttt{lake env lean} em projeto pinado (Lean v4.34.1, Mathlib rev.\
v4.34.1, 4760 oleans pré-compilados oficiais); auditoria \texttt{\#print
axioms} por teorema; bloqueio \textit{fail-closed} para \texttt{sorry},
\texttt{admit} e axiomas não declarados. \textit{Bancada.} Piscina de 40
agentes sintéticos por semente com qualidade latente $q$ oculta do roteador;
oráculo $cobertura\times(0{,}4+0{,}6q)$; 5 sementes $\times$ 40 tarefas = 200
decisões; estratégias: roteador real, aleatório, menor carga, maior cobertura;
métricas: arrependimento médio e taxa de acerto, IC95 por aproximação normal.
\textit{Ablações.} Confiança embaralhada, anticorrelacionada e ruidosa; média
igualitária das cabeças; remoção da cabeça de confiança; análise pareada e
sensibilidade do oráculo em $q_w\in\{0{,}3;0{,}6;0{,}9\}$.
\textit{Contraprova numérica.} 20\,000 vetores aleatórios (magnitudes até
$\pm1000$) contra a implementação real em ponto flutuante.
\textit{Revisões.} Duas revisões independentes internas (bancada e teoremas)
via executores externos; revisores por agente indisponíveis no ambiente
(falha registrada). Ameaças: oráculo sintético e arbitrário; confiança
correlacionada com $q$ (proxy privilegiada); baselines ingênuos; 5 sementes;
lacuna floats versus reais; sem tarefas reais nem concorrentes executados.

\section{Resultados}

Os 11 lemas listados na Tabela~\ref{tab:ax} obtiveram \textit{machine-checked} (saída 0 do kernel em 2026-10-10),
com axiomas restritos a \textit{propext}, \textit{Classical.choice} e
\textit{Quot.sound} (Tabela~\ref{tab:ax}).
\begin{table}[h]
\centering
\caption{Auditoria de axiomas por teorema (\texttt{\#print axioms}).}
\label{tab:ax}
\begin{tabular}{ll}
\toprule
Teorema & Axiomas \\
\midrule
\texttt{softmax\_soma\_um} & propext, Classical.choice, Quot.sound \\
\texttt{combinacao\_convexa\_limitada} & propext, Classical.choice, Quot.sound \\
\texttt{roteamento\_guloso\_otimo} & propext, Classical.choice, Quot.sound \\
\texttt{softmax\_invariante\_max} & propext, Classical.choice, Quot.sound \\
\texttt{softmax\_mem\_Ioo} & propext, Classical.choice, Quot.sound \\
\texttt{combinacao\_convexa\_limitada\_baixo} & propext, Classical.choice, Quot.sound \\
\texttt{softmax\_peso\_total\_singleton} & propext, Classical.choice, Quot.sound \\
\texttt{peso\_cresce\_com\_escore} & propext, Classical.choice, Quot.sound \\
\texttt{topk\_dominam\_resto} & propext, Classical.choice, Quot.sound \\
\texttt{temperatura\_aguca\_maximo} & propext, Classical.choice, Quot.sound \\
\texttt{sem\_overflow\_apos\_max} & propext, Classical.choice, Quot.sound \\
\bottomrule
\end{tabular}
\end{table}
A primeira versão de T7 continha \texttt{sorryAx} por erro de aplicação; a
auditoria detectou e a correção foi re-verificada --- evidência de que a
trilha funciona. Na bancada (Tabela~\ref{tab:bench}), o roteador obteve
arrependimento $0{,}042\pm0{,}011$ e acerto $0{,}625\pm0{,}067$, com ICs sem
sobreposição com os baselines ($0{,}19$--$0{,}22$); latência média $2{,}17$\,ms;
soma do softmax igual a 1 até $1{,}1\times10^{-16}$.
\begin{table}[h]
\centering
\caption{Bancada principal (200 decisões).}
\label{tab:bench}
\begin{tabular}{lcc}
\toprule
Estratégia & Arrependimento & Acerto \\
\midrule
Atenção (roteador real) & $0{,}042\pm0{,}011$ & $0{,}625\pm0{,}067$ \\
Aleatório & $0{,}188\pm0{,}022$ & $0{,}200\pm0{,}055$ \\
Menor carga & $0{,}217\pm0{,}025$ & $0{,}215\pm0{,}057$ \\
Maior cobertura & $0{,}212\pm0{,}022$ & $0{,}165\pm0{,}051$ \\
\bottomrule
\end{tabular}
\end{table}
A contraprova numérica registrou 0 violações em 20\,000 vetores (pior desvio
de soma $2{,}2\times10^{-16}$; invariância a deslocamento exata). Nas
ablações, a atenção venceu em todos os mundos e a fusão ponderada superou
pesos iguais (vitória em $0{,}99$ das decisões); registra-se o achado
desconfortável de que, sem a cabeça de confiança, o roteador equivale ao
acaso neste simulador --- a vantagem depende de reputação calibrada.
Bateria de regressão: 63 testes aprovados; diagnóstico do ambiente sem
falhas (18/20 verificações, 2 advertências pré-existentes de CLIs externas).

\section{Impactos teóricos, práticos e científicos}

\textit{Teórico.} O núcleo matemático do roteamento deixa de ser heurística
opaca: cada propriedade operacional (pesos que somam 1, decisão invariante a
deslocamento, escores limitados, seleção máxima) corresponde a um enunciado
verificado por máquina com axiomas auditados. A invariância (T4) fundamenta a
prática de estabilidade numérica; o par (T5, T7) elimina o caso omisso do
singleton; o par (T2, T6) dá conteúdo formal à exigência de soma unitária.
\textit{Prático.} As decisões do roteador passam a carregar pesos calibrados
(soma 1 até $10^{-16}$) e regra de seleção provadamente máxima quanto à
utilidade definida, ao custo medido de $2{,}17$\,ms por decisão; a seleção
parcial segue como melhoria medida pendente. \textit{Científico.} A trilha
funcionou como instrumento: a auditoria de axiomas capturou um
\texttt{sorryAx} silencioso, a auditoria de referências descartou uma citação
suposta, e as ablações quantificaram a dependência da cabeça de confiança ---
achados que integram o relato em vez de escondê-lo. \textit{Limites do
escopo:} superioridade sobre plataformas existentes não foi testada; a
avaliação restringiu-se ao simulador descrito no Método; propriedades de
Transformers treinados não foram avaliadas.

\section{Discussão}

As revisões internas classificaram a evidência como preliminar restrita ao
simulador (bancada) e os teoremas como revisão maior: os lemas são
elementares, a analogia com o Transformer permanece cosmética sem
correspondência consulta--chave--valor formalizada, e persiste a lacuna entre
reais formalizados e floats executados, top-\textit{k}, gates e empates. A
renúncia a superioridade universal delimita, mas não prova, eficiência
operacional. Originalidade, comparação \textit{head-to-head} com plataformas
maduras e reprodução independente permanecem como trabalho futuro
indispensável antes de qualquer alegação de inovação ou disrupção.

\section{Conclusão}

Entrega-se um núcleo matemático do roteamento por atenção verificado por
máquina, calibrado internamente contra baselines ingênuos e auditado até os
axiomas, com trilha de reprodução integralmente publicada (código, sementes,
versões, hashes e comandos no Apêndice). Reivindica-se exclusivamente
otimalidade relativa e evidência preliminar.

\begin{thebibliography}{9}
\bibitem{moura2021} MOURA, L. de; ULLRICH, S. The Lean 4 theorem prover and
programming language. \textit{Lecture Notes in Computer Science}, p. 625--635,
2021. DOI: \url{https://doi.org/10.1007/978-3-030-79876-5_37}.
Verificada na fonte primária (Crossref).
\bibitem{mathlib2020} THE MATHLIB COMMUNITY. The lean mathematical library.
In: \textit{Proceedings of the 9th ACM SIGPLAN International Conference on
Certified Programs and Proofs}, 2020. p. 367--381. DOI:
\url{https://doi.org/10.1145/3372885.3373824}.
Verificada na fonte primária (Crossref).
\bibitem{vaswani2017} VASWANI, A. et al. Attention is all you need. In:
\textit{Advances in Neural Information Processing Systems 30}, Long Beach,
2017. p. 5998--6008. (Sem DOI; obra de conferência amplamente referenciada.)
\end{thebibliography}

\appendix
\sloppy
\section{Trilha de reprodução}

Lean 4 v4.34.1, Lake 5.0.0, elan 4.2.4; projeto \texttt{lean-mathlib-demo},
\texttt{lean-toolchain} \texttt{leanprover/lean4:v4.34.1}, Mathlib rev.\
\texttt{v4.34.1}; commits verificados por \texttt{git ls-remote} em
2026-10-09: lean4 \texttt{5e96ee11\ldots}, mathlib4 \texttt{a02a63c2\ldots},
batteries \texttt{b71aaf60\ldots}, aesop \texttt{04879b1e\ldots},
z3 \texttt{0d0e28c3\ldots}, sympy \texttt{04a7da3a\ldots}; licenças Apache-2.0
(Lean), MIT (z3), BSD (sympy) com SHA-256 custodiados. Sementes 11, 22, 33,
44, 55; 40 tarefas por semente. Arquivos: teoremas em
\url{plugins/lean-nucleo-verificacao/teoremas/}; bancada, ablações,
contrafuzz e auditoria de axiomas em
\url{plugins/lean-nucleo-verificacao/bancada/}; testes
\texttt{test\_r747}, \texttt{test\_r754}, \texttt{test\_r755}; especificações
SPEC-935-R747/R754/R755. Reprodução: \texttt{lake env lean
RoteamentoAtencao.lean} (saída 0) e \texttt{pytest} das 9 suítes (63
aprovados). Depósito em repositório aberto (OSF): pendente de ação do autor;
materiais sob custódia nos caminhos citados. Contatos e ORCID no cabeçalho.

\end{document}
